\documentclass[journal, a4paper]{IEEEtran}

\usepackage[T1]{fontenc}
\usepackage[utf8]{inputenc}
\usepackage[spanish, es-noshorthands]{babel}
\usepackage{amsmath, amssymb, bm}
\usepackage{graphicx}
\usepackage{booktabs}
\usepackage{tabularx}
\usepackage{ragged2e}
\usepackage[american]{circuitikz}
\usepackage{pgfplots}
\pgfplotsset{compat=1.18}
\usepackage{hyperref}

\renewcommand{\tablename}{Tabla}

\begin{document}

% ==========================================
% TÍTULO Y AUTORES
% ==========================================
% INDIQUE el problema técnico, no solo "Informe de Laboratorio".
% Ejemplo: "Verificación Experimental de [Fenómeno/Teorema] en Redes Resistivas"
\title{Título del Informe: [Fenómeno/Teorema] en Redes Resistivas}

\author{
    Nombre del Estudiante 1 (ID: XXXXXXXX), Nombre del Estudiante 2 (ID: XXXXXXXX)\\
    Ingeniería Mecatrónica, Universidad Católica Boliviana ``San Pablo'' -- Sede Tarija\\
    IMT-121 --- Circuitos Electrónicos I\\
    Fecha: \today \quad | \quad Grupo: [Número o Identificador]
}

\maketitle

% ==========================================
% ABSTRACT Y PALABRAS CLAVE
% ==========================================
\begin{abstract}
% STATE what you tested.
% SHOW the analytical prediction.
% REPORT the simulation and measurement.
% QUANTIFY the difference.
% EXPLAIN what the comparison means.
% Escriba entre 80 y 120 palabras en un solo párrafo. No incluya detalles del procedimiento.
En este documento se presenta la verificación experimental de...
\end{abstract}

\begin{IEEEkeywords}
Término 1, Término 2, Término 3. % Ej: Superposición, Equivalente de Thévenin, LTSpice.
\end{IEEEkeywords}

% ==========================================
% I. INTRODUCCIÓN Y OBJETIVOS
% ==========================================
\section{Introducción y Objetivo}
% Defina la pregunta de ingeniería. ¿Qué relación o teorema se está evaluando?
% Termine esta sección (1 a 3 párrafos cortos) con un objetivo explícito.
Este experimento investiga si una red resistiva lineal puede...

% ==========================================
% II. MODELO DEL CIRCUITO Y PREDICCIONES
% ==========================================
\section{Modelo del Circuito y Predicciones Analíticas}
% Presente el modelo analítico antes que las mediciones experimentales.

\subsection{Definición del Circuito y Referencias}
% Incluya un esquemático con las referencias de voltaje, corriente, polaridades, nodos y el puerto a-b.
\begin{figure}[htbp]
    \centering
    % Placeholder para el circuito (usar circuitikz)
    \begin{circuitikz}[scale=0.8, transform shape]
        % Dibuje aquí el circuito
        \draw (0,0) to[R, l={$R$}] (2,0);
    \end{circuitikz}
    \caption{Esquemático del circuito bajo prueba indicando referencias y puerto de análisis.}
    \label{fig:circuito}
\end{figure}

\subsection{Relaciones Fundamentales}
% Incluya SOLO las ecuaciones estrictamente necesarias para el experimento.
% Muestre suficiente desarrollo intermedio para que el método se entienda. No provea páginas de álgebra.
Las ecuaciones que gobiernan el comportamiento del circuito son:
\begin{equation}
    % Ecuación de ejemplo
    V_{th} = V_{oc}
\end{equation}

\subsection{Predicciones Analíticas}
% Resuma en una tabla compacta los valores esperados antes de la simulación y medición.
\begin{table}[htbp]
    \centering
    \caption{Predicciones Analíticas Previas}
    \label{tab:predicciones}
    \begin{tabularx}{\columnwidth}{>{\RaggedRight\arraybackslash}X >{\centering\arraybackslash}X c}
        \toprule
        \textbf{Magnitud} & \textbf{Predicción Analítica} & \textbf{Unidad} \\
        \midrule
        Variable 1 & Valor & Unidad \\
        Variable 2 & Valor & Unidad \\
        \bottomrule
    \end{tabularx}
\end{table}

% ==========================================
% III. METODOLOGÍA
% ==========================================
\section{Metodología}

\subsection{Simulación (SPICE)}
% Describa el software usado, la configuración del circuito, qué parámetros se barrieron (si aplica) y qué cantidades se extrajeron.

\subsection{Montaje Experimental}
% Documente la implementación en protoboard, configuración de fuentes de alimentación, instrumentos usados, y punto/nodo de referencia (tierra).
% No invente números de modelo de instrumentos si no los registró.

\subsection{Procedimiento}
% Use una secuencia numerada y concisa indicando qué cambió, qué permaneció fijo, qué se midió y cómo se construyó la comparación.
\begin{enumerate}
    \item Se configuró...
    \item Se mantuvo fijo...
    \item Se midió...
\end{enumerate}

% ==========================================
% IV. RESULTADOS
% ==========================================
\section{Resultados}
% Reporte la evidencia sin mezclarla inmediatamente con largas interpretaciones.

% Utilice la siguiente fórmula para el error relativo cuando sea apropiado (no cuando el teórico sea cero):
% \varepsilon_r = \frac{|x_{\text{medido}} - x_{\text{teórico}}|}{|x_{\text{teórico}}|} \times 100\%

\begin{table*}[htbp]
    \centering
    \caption{Comparación de Resultados: Teoría, Simulación y Medición}
    \label{tab:resultados}
    \renewcommand{\arraystretch}{1.2}
    \begin{tabularx}{\textwidth}{>{\RaggedRight\arraybackslash}X *{4}{>{\centering\arraybackslash}X}}
        \toprule
        \textbf{Magnitud} & \textbf{Teoría} & \textbf{SPICE} & \textbf{Medición} & \textbf{Error Relativo (\%)} \\
        \midrule
        Variable 1 & Valor & Valor & Valor & Valor \\
        Variable 2 & Valor & Valor & Valor & Valor \\
        \bottomrule
    \end{tabularx}
\end{table*}

% Agregue aquí la figura principal del experimento (ej. gráfica P_L vs R_L).
% Evite llenar el reporte con capturas de pantalla de osciloscopios o simuladores si no aportan valor.

% ==========================================
% V. DISCUSIÓN
% ==========================================
\section{Discusión}
% Esta es la sección principal de razonamiento de ingeniería.
% Responda explícitamente:
% 1. ¿Acordaron razonablemente la teoría, SPICE y la medición?
% 2. ¿Cuál fue la magnitud de la discrepancia más importante?
% 3. ¿Qué causas físicas o procedimentales plausibles podrían explicarla? (ej. tolerancia de resistores, resistencia de cables, tierra incorrecta).
% MODELO A SEGUIR: Observación -> Hipótesis Técnica -> Evidencia.
% NO escriba simplemente "error experimental o humano".

% ==========================================
% VI. CONCLUSIONES
% ==========================================
\section{Conclusiones}
% Proporcione de 3 a 5 conclusiones técnicas concisas ligadas a su evidencia.
% NO escriba "los objetivos se cumplieron".
\begin{itemize}
    \item La medición de la tensión de circuito abierto difirió del $V_{th}$ analítico en un X\%, lo cual...
    \item La red cargada y su equivalente de Thévenin produjeron tensiones terminales similares, soportando...
\end{itemize}

% ==========================================
% REFERENCIAS
% ==========================================
\begin{thebibliography}{00}
\bibitem{ref1}
% Use formato IEEE real para la guía de lab, el libro de texto o documentación de LTSpice.
% Ejemplo de formato: Autor, \emph{Título del Documento}, Edición. Ciudad: Editorial, Año.
B. Quiroga Turdera, \emph{Guía de Laboratorio IMT-121}, Universidad Católica Boliviana, 2026.
\end{thebibliography}

% ==========================================
% APÉNDICE (Opcional)
% ==========================================
% \section*{Apéndice}
% Úselo SOLO para tablas raw muy largas, cálculos extensos o código. El reporte principal debe entenderse sin el apéndice.

\end{document}
